\documentclass[11pt]{article}
\usepackage{amsmath,amssymb}
\usepackage[margin=1in]{geometry}
\usepackage{booktabs}

\newcommand{\ci}[4]{(#1#2\,|\,#3#4)}          % chemist (mu nu | la si)
\newcommand{\pd}[1]{\partial_{#1}}

\title{Permutational symmetry of two-electron derivative integrals\\
       \large (chemist notation, no antisymmetrization)}
\author{}
\date{}

\begin{document}
\maketitle
\vspace{-3em}

\section{Setup}

Work with the ordinary (non-antisymmetrized) two-electron repulsion integral in
chemist notation, real orbitals,
%
\begin{equation}
\ci{\mu}{\nu}{\lambda}{\sigma}
= \int\!\!\int
   \chi_\mu(\mathbf r_1)\chi_\nu(\mathbf r_1)\,\frac{1}{r_{12}}\,
   \chi_\lambda(\mathbf r_2)\chi_\sigma(\mathbf r_2)\;
   d\mathbf r_1\,d\mathbf r_2 .
\end{equation}
%
It carries an order-8 permutation group generated by
%
\begin{equation}
\ci{\mu}{\nu}{\lambda}{\sigma} = \ci{\nu}{\mu}{\lambda}{\sigma}
= \ci{\mu}{\nu}{\sigma}{\lambda}
= \ci{\lambda}{\sigma}{\mu}{\nu} ,
\label{eq:gens}
\end{equation}
%
i.e.\ swap the two bra functions, swap the two ket functions, and swap the bra
charge distribution $\mu\nu$ (electron 1) with the ket charge distribution
$\lambda\sigma$ (electron 2). The last of these --- the bra--ket swap --- is the
one that matters below.

\section{The symmetry is the same at every order}

Each relation in \eqref{eq:gens} is an identity in the function labels that holds
at \emph{every} value of any perturbation parameters $\boldsymbol\lambda$ on which
the integral depends. A permutation $\pi$ of $\{\mu,\nu,\lambda,\sigma\}$ is
independent of $\boldsymbol\lambda$; a derivative $\pd{X},\ \pd{X}\pd{Y}$ acts on
$\boldsymbol\lambda$. They commute, so differentiating an identity valid for all
$\boldsymbol\lambda$ gives the same identity among the derivatives:
%
\begin{equation}
\ci{\mu}{\nu}{\lambda}{\sigma}=\ci{\lambda}{\sigma}{\mu}{\nu}
\ \ \forall\boldsymbol\lambda
\;\Longrightarrow\;
\ci{\mu}{\nu}{\lambda}{\sigma}^{(X)}=\ci{\lambda}{\sigma}{\mu}{\nu}^{(X)},\quad
\ci{\mu}{\nu}{\lambda}{\sigma}^{(XY)}=\ci{\lambda}{\sigma}{\mu}{\nu}^{(XY)},
\end{equation}
%
and likewise for the other two generators.

\begin{quote}
\textbf{The exact derivative integral carries the full order-8 group at every
order, and for any choice of perturbation parameters --- same type or different
type.} There is no reduction of symmetry with order or with the kind of
perturbation.
\end{quote}
%
A derivative integral is nonzero only where the integral actually depends on the
parameter. The two-electron integral depends on the nuclear positions (through the
orbital centers) but not on an external electric field, so field derivatives of
$\ci{\mu}{\nu}{\lambda}{\sigma}$ vanish identically and only nuclear derivatives
are nonzero. This governs \emph{whether} a derivative integral exists, not its
symmetry: whenever one is nonzero, it has the full group.

\section{What \texttt{mo\_tei\_deriv2} actually returns}

\paragraph{The target.} Let $X$ and $Y$ be any two perturbation coordinates and
write $\chi^{X}\equiv\pd{X}\chi$. By the product rule the complete mixed second
derivative distributes $\pd{X}$ over the four functions and $\pd{Y}$ over the four
functions, giving $16$ terms:
%
\begin{align}
(\mu\nu|\lambda\sigma)^{(XY)}
={}& (\mu^{XY}\nu|\lambda\sigma) + (\mu^{X}\nu^{Y}|\lambda\sigma)
   + (\mu^{X}\nu|\lambda^{Y}\sigma) + (\mu^{X}\nu|\lambda\sigma^{Y}) \notag\\
 +{}& (\mu^{Y}\nu^{X}|\lambda\sigma) + (\mu\nu^{XY}|\lambda\sigma)
   + (\mu\nu^{X}|\lambda^{Y}\sigma) + (\mu\nu^{X}|\lambda\sigma^{Y}) \notag\\
 +{}& (\mu^{Y}\nu|\lambda^{X}\sigma) + (\mu\nu^{Y}|\lambda^{X}\sigma)
   + (\mu\nu|\lambda^{XY}\sigma) + (\mu\nu|\lambda^{X}\sigma^{Y}) \notag\\
 +{}& (\mu^{Y}\nu|\lambda\sigma^{X}) + (\mu\nu^{Y}|\lambda\sigma^{X})
   + (\mu\nu|\lambda^{Y}\sigma^{X}) + (\mu\nu|\lambda\sigma^{XY}) .
\label{eq:complete}
\end{align}
%
Row $k$ collects the four terms with $\pd{X}$ on the $k$-th function
($\mu,\nu,\lambda,\sigma$ respectively) and $\pd{Y}$ ranging over all four. The
four diagonal terms $(\,\cdots^{XY}\cdots)$ --- both derivatives on one function
--- are nonzero only when $X$ and $Y$ are coordinates of the \emph{same} atom; if
$X$ is a coordinate of atom $A$ and $Y$ of a different atom $B$ they vanish (an
orbital sits on one atom and is unmoved by displacing another), leaving the $12$
off-diagonal terms. By \S2 this complete derivative obeys the full group
\eqref{eq:gens}; in particular the bra--ket relation
$(\mu\nu|\lambda\sigma)^{(XY)}=(\lambda\sigma|\mu\nu)^{(XY)}$.

\paragraph{What Psi4 returns.} For a shell quartet the integral engine
(\texttt{compute\_shell\_deriv2}) supplies the second derivatives of
$(\mu\nu|\lambda\sigma)$ with respect to the four orbital centers, which we label
$A\!=\!\mu$, $B\!=\!\nu$ (bra) and $C\!=\!\lambda$, $D\!=\!\sigma$ (ket). For a
requested atom pair $(\mathrm{atom}_1,\mathrm{atom}_2)$ with
$\mathrm{atom}_1\neq\mathrm{atom}_2$, \texttt{ao\_tei\_deriv2} /
\texttt{mo\_tei\_deriv2}:
%
\begin{enumerate}
\item enumerate the six off-diagonal center-pairs in one fixed order,
$AB,\,AC,\,AD,\,BC,\,BD,\,CD$ (the upper triangle; the reversed pairs
$BA,\,CA,\dots$ are never formed);
\item keep those whose (first center, second center) lie on
$(\mathrm{atom}_1,\mathrm{atom}_2)$ in that order --- i.e.\ only when the
$\mathrm{atom}_1$-center precedes the $\mathrm{atom}_2$-center in the
$\mu<\nu<\lambda<\sigma$ ordering; and
\item multiply each kept contribution by $2$.
\end{enumerate}
%
The returned array is therefore
%
\begin{equation}
\texttt{deriv2}(\mathrm{atom}_1,\mathrm{atom}_2)
= 2\!\!\!\sum_{\substack{(c_i,c_j)\ \text{canonical},\ i<j\\
   c_i\ \text{on}\ \mathrm{atom}_1,\ c_j\ \text{on}\ \mathrm{atom}_2}}\!\!\!
   \frac{\partial^2 (\mu\nu|\lambda\sigma)}{\partial c_i\,\partial c_j}\, ,
\label{eq:raw}
\end{equation}
%
twice the upper-triangular (canonically ordered) center-pairs and none of the
lower-triangular ones. In terms of \eqref{eq:complete} it keeps the terms whose
$X$-differentiated function precedes its $Y$-differentiated function in the
$\mu\nu\lambda\sigma$ ordering, doubled, and drops the rest. It therefore violates
the bra--ket relation $(\mu\nu|\lambda\sigma)^{(XY)}=(\lambda\sigma|\mu\nu)^{(XY)}$
by an $O(1)$ amount.

\paragraph{Why the doubling.} It is exact for the intended use. The molecular
Hessian contracts these integrals with the bra--ket-symmetric two-particle
density, and against a symmetric density $2\times(\text{upper triangle})$ equals
$(\text{upper}+\text{lower})$, the complete derivative. The engine thus forms half
the center-pairs and doubles. The returned tensor is correct under a
symmetric-density contraction but is not itself the symmetric integral.

\section{Recovering the complete integral}

Used outside a symmetric-density contraction, the true symmetric derivative must
be reconstructed. The bra--ket transpose of \eqref{eq:raw} supplies exactly the
lower-triangular center-pairs the enumeration skipped,
$\texttt{deriv2}^{\,\mu\nu\leftrightarrow\lambda\sigma}=2\times(\text{lower
triangle})$, so the bra--ket average
%
\begin{equation}
\tfrac12\Bigl(\texttt{deriv2}
   + \texttt{deriv2}^{\,\mu\nu\leftrightarrow\lambda\sigma}\Bigr)
= (\mu\nu|\lambda\sigma)^{(XY)}
\label{eq:avg}
\end{equation}
%
restores the missing pairs and cancels the factor $2$ (confirmed by finite
difference to $\sim\!10^{-7}$). This is the symmetrization the derivative codes
apply:
%
\begin{itemize}
\item spin-orbital path: \texttt{so\_eri2} applies
$\tfrac12(\texttt{ch}+\texttt{ch}^{\,\mu\nu\leftrightarrow\lambda\sigma})$ to the
chemist array \emph{before} any conversion or spin bookkeeping, so its consumers
(\texttt{cphf}, \texttt{CorrelatedDerivs}) receive a completed integral;
\item spatial path: the raw chemist array is passed through unaveraged and
\texttt{cphf} applies the same average where it assembles the nuclear--nuclear
skeletons.
\end{itemize}
%
It is required only for the nuclear--nuclear second derivative, the only place a
nonzero second-derivative two-electron integral is formed. First derivatives are
returned complete and need no averaging.

\section{Transforming \texttt{ao\_tei\_deriv2} directly}

\texttt{mo\_tei\_deriv2} is \texttt{mo\_eri\_helper} applied to the
\texttt{ao\_tei\_deriv2} blocks, and that helper does not transform the four
indices in place. To keep its \texttt{DGEMM}s contiguous it contracts the third MO
coefficient against the AO $\sigma$ and the fourth against the AO $\lambda$ (via
two explicit buffer reorderings), so its output is
%
\begin{equation}
\texttt{mo\_tei\_deriv2}[p,q,r,s]
  = \sum_{\mu\nu\lambda\sigma} C_{\mu p}\,C_{\nu q}\,C_{\sigma r}\,C_{\lambda s}\,
    (\mu\nu|\lambda\sigma)^{(XY)} ,
\label{eq:ketreorder}
\end{equation}
%
i.e.\ the two \emph{ket} MO indices $r,s$ ride on the AO indices $\sigma,\lambda$
--- the ket pair is transposed relative to a straight index-by-index transform,
which would put $r$ on $\lambda$ and $s$ on $\sigma$.

A consumer that transforms \texttt{ao\_tei\_deriv2} \emph{directly}, bypassing
\texttt{mo\_tei\_deriv2} --- as the memory-lean AO Hessian-skeleton route does
(\texttt{Derivatives.eri2\_mo\_component} / \texttt{so\_eri2\_mo\_component}, which
hold the nine AO blocks for a pair and transform one Cartesian component at a
time) --- must therefore swap the AO ket pair (\texttt{transpose(0,1,3,2)}) before
the transform to reproduce \eqref{eq:ketreorder}.

This is not cosmetic. The bra--ket average \eqref{eq:avg} completes the raw
\emph{in the ket order Psi4 emits}, and the raw is not ket-symmetric: the
enumeration \eqref{eq:raw} is keyed on the global $\mu<\nu<\lambda<\sigma$
ordering, which $\lambda\leftrightarrow\sigma$ breaks (it sends the
upper-triangular center-pair $CD$ to the lower-triangular $DC$). Averaging a
ket-swapped raw therefore yields a \emph{different}, ket-permuted tensor, not the
symmetric integral. The difference cancels only once the block is contracted with
a ket-symmetric two-particle density: CISD and MP2 symmetrize their 2-PDM and are
unaffected, but the CCSD (and CCSD(T)) cumulant $\Gamma$ is not ket-symmetric, so
omitting the reorder corrupts the spatial CCSD/CCSD(T) Hessian by $O(10^{-2})$
while CISD/MP2 stay correct.

\section{Summary}

\begin{itemize}
\item Exact two-electron derivative integrals have the full order-8 permutational
symmetry at every order and for every combination of perturbation parameters.
\item For $\mathrm{atom}_1\neq\mathrm{atom}_2$, \texttt{mo\_tei\_deriv2} returns
twice the upper-triangular (canonically ordered) center-pair second derivatives,
\eqref{eq:raw}. The doubling makes it correct against a symmetric-density Hessian
contraction, but leaves it bra--ket asymmetric and not equal to the integral.
\item The complete symmetric integral is the bra--ket average
$\tfrac12(\texttt{deriv2}+\texttt{deriv2}^{\,\mu\nu\leftrightarrow\lambda\sigma})$,
\eqref{eq:avg}. \texttt{so\_eri2} applies it internally; \texttt{cphf} applies it
for the spatial skeletons. First derivatives need none.
\item \texttt{mo\_eri\_helper} transposes the ket pair while transforming
\eqref{eq:ketreorder}, so a direct \texttt{ao\_tei\_deriv2} consumer must
\texttt{transpose(0,1,3,2)} the AO block to match. The raw is not ket-symmetric,
so skipping this is silent against a ket-symmetric density (CISD/MP2) but wrong
against the CCSD/CCSD(T) $\Gamma$.
\end{itemize}

\end{document}
